\chapter{Iteration 2: Design and Implementation}
\label{chap:iter2}
\textbf{Baseline.} This chapter covers work from 13~October to the end of FYP-1
on \texttt{develop} (commit \texttt{ce49dcc}, \tbd{update to the submission tag}).

\section{Iteration Objectives}
Iteration 2 had three objectives, taken from the Iteration~1 outcome and the
proposal feedback: (1) join live capture to memory, (2) make the memory bounded
and persistent across sessions (FR3.5, FR4), and (3) give the patient and
caregiver a usable product (proposal action 4). A fourth, unplanned objective
was to replace the detector after the benchmark of Section~\ref{sec:experiments}.

\section{Selected Use Cases / Features}
UC-01 (live), UC-03, UC-04, UC-06, UC-12 (compaction) and UC-13 (rooms) were
selected, plus UC-02 by speech, and the new UC-14 and UC-15.
Table~\ref{tab:i2-features} lists the features delivered.

\begin{table}[H]\centering\footnotesize
\caption{Iteration~2 features and their status.}\label{tab:i2-features}
\begin{tabularx}{\textwidth}{@{}Y L{3.6cm} L{2.0cm} l@{}}\toprule
\rowcolor{hdrfill}\textbf{Feature} & \textbf{Main files} & \textbf{FRs} & \textbf{Status}\\\midrule
Live capture contracts, keyframes, gaps, buffer, crash resume & \path{ingestion/live_pipeline.py}, \texttt{keyframe.py}, \texttt{buffer.py} & FR1, FR2 & Implemented\\
Promotion into segments and Modal runner & \texttt{live\_handoff.py}, \texttt{segment\_runner.py} & FR2.2 & Implemented\\
Gravity frame (RANSAC floor fit) & \path{geometry/gravity.py} & FR3.8 & Implemented\\
Nine relations in the gravity frame & \path{memory/relations.py} & FR3.8 & Implemented\\
State runs and events (schema 2) & \path{memory/runs.py}, \texttt{scene\_graph.py} & FR3.2, FR3.3 & Implemented\\
SQLite home store and cross-session matching & \texttt{store.py}, \texttt{matching.py}, \texttt{ingest.py} & FR3.5, FR3.9 & Implemented\\
Rooms, activities, room registry & \texttt{rooms.py}, \texttt{activities.py}, \texttt{home\_registry.py} & FR3.4, FR3.6 & Partial\\
Compaction with verification and roll-up & \texttt{compaction.py}, \texttt{summaries.py} & FR4 & Implemented\\
Evidence crops & \texttt{evidence.py} & FR7.2 & Implemented\\
SAM~3 detector, DINOv3 re-ID & \path{perception/detector.py} & FR3.1 & Implemented\\
Ranked candidates (tiered protocol) & \texttt{reasoning.py} & FR5.8 & Implemented\\
NLQ engine, dictionary (525 items), \path{/ask} & \path{nlq/engine.py}, \texttt{dictionary.py} & FR5.1, FR5.5 & Partial\\
Modal-served API & \path{remote/serve.py} & -- & Implemented\\
Patient App and Caregiver App & \texttt{apps/Patient App}, \texttt{apps/Caregiver App} & FR6, FR7, FR8.6, FR8.7 & Implemented (no auth)\\
Evaluation suite and EgoLife study & \path{evaluation/} & NFR-P & Implemented\\
\bottomrule\end{tabularx}\end{table}

\section{Design}
Figures~\ref{fig:i2-graph}--\ref{fig:i2-query} are the Iteration~2 class
diagrams, drawn from the code. Only fields and methods that matter to the design
are shown.

\aegisfig{ch7-class-memory-graph}{Memory graph classes (schema 2).}{fig:i2-graph}
\aegisfig{ch7-class-memory-runs}{Observations, state runs and events.}{fig:i2-runs}
\aegisfig{ch7-class-memory-store}{\texttt{HomeStore} and its records.}{fig:i2-store}
\aegisfig{ch7-class-memory-ingest}{Session ingest, matching and compaction.}{fig:i2-ingest}
\aegisfig{ch7-class-api-apps}{API models and the two Flutter apps.}{fig:i2-apps}
\aegisfig{ch7-class-query}{Query: ranking, NLQ engine and dictionary.}{fig:i2-query}

\clearpage
\section{Implementation}
This section describes the four algorithms added in Iteration~2.

\subsection{Relation classification}
\label{sec:relations-algo}
Each object box is compared with every anchor box in the gravity frame
(Figure~\ref{fig:flow-relations}). If gravity confidence is below 0.5 the
relation is \texttt{near} and marked uncertain. Otherwise, with a vertical
tolerance of 0.15 and a horizontal overlap threshold of 0.5, the classifier
decides in order: \texttt{carried} (a hand cue), \texttt{inside} (container
anchor, box within it), \texttt{on\_surface}, \texttt{under} (anchor open
underneath), \texttt{above}, \texttt{on\_floor} (within 6\% of the scene height
of the floor), \texttt{next\_to}, \texttt{near}, and finally \texttt{in\_room}.
Event types are \texttt{placed}, \texttt{picked\_up}, \texttt{moved},
\texttt{relation\_changed}, \texttt{room\_changed} and
\texttt{moved\_while\_unobserved}.

\aegisfigtall{ch7-flow-relations}{Relation classification flowchart.}{fig:flow-relations}

\subsection{State runs}
Observations of one instance are sorted by time. A run starts only when two
consecutive observations agree on anchor and relation (or one final observation
with confidence $\geq$ 0.5); instances with fewer than two observations are
dropped. Each change between runs gives one event. An object seen in the same
place for an hour is therefore one run, not 7{,}200 timeline rows.

\subsection{Cross-session matching}
\label{sec:matching-algo}
A new session is matched to the home in three levels
(Figure~\ref{fig:flow-matching}). \textbf{Rooms} match by landmark-signature
similarity (at least 0.5, a margin of 0.05 over the runner-up, at least one shared
landmark). \textbf{Anchors} are aligned by a rigid fit whose RMS residual must
stay below 25\% of the scene spread. \textbf{Objects} match by DINOv3 appearance
distance below 0.35 with a 0.1 margin; an object found at a different anchor
than its last-known state records \texttt{moved\_while\_unobserved}.

\aegisfig{ch7-flow-matching}{Cross-session matching flowchart.}{fig:flow-matching}

\subsection{Keyframe selection}
The live sampler keeps a frame (Figure~\ref{fig:flow-keyframe}) only if its mean
brightness is between 20 and 245, its Laplacian variance is at least 15 and its
change from the last kept frame (on a 64$\times$36 thumbnail) is at least 0.008;
a frame is forced after 2\,s without a keyframe so a still scene is still
recorded.

\aegisfig{ch7-flow-keyframe}{Keyframe selection flowchart.}{fig:flow-keyframe}

\subsection{Compaction}
Each local day (Asia/Karachi) older than three days is summarised by templates,
verified and only then deleted; important events stay seven days. Day summaries
older than 30 days roll up into weeks, and weeks older than 365 days into months.
The last-known run and evidence crop of every object are kept.

\section{Technical Implementation Details}
Table~\ref{tab:params} lists the defaults in \path{core/config.py}. All are
overridable from \path{configs/*.yaml}.

\begin{table}[H]\centering\footnotesize
\caption{Main Iteration~2 parameters.}\label{tab:params}
\begin{tabularx}{\textwidth}{@{}L{4.6cm} l Y@{}}\toprule
\rowcolor{hdrfill}\textbf{Parameter} & \textbf{Default} & \textbf{Effect}\\\midrule
\texttt{live.sample\_fps} & 2.0 & Frames sampled from MJPEG\\
\texttt{live.max\_duration\_seconds} & 600 & Session limit\\
\texttt{live.reconnect\_timeout\_seconds} & 30 & End after this long without frames\\
\texttt{live.buffer\_capacity} & 30 & Frames held during a stall\\
\texttt{live.segment\_frames} & 120 & Frames per promoted segment\\
\texttt{max\_gap\_seconds} & 30 & Gap that splits a segment\\
\texttt{memory.anchor\_cluster\_radius} & 0.40 & Anchor clustering (VGGT units)\\
\texttt{memory.association\_radius} & 0.45 & Same-instance radius\\
\texttt{memory.appearance\_threshold} & 0.35 & DINOv3 distance for re-ID\\
\texttt{memory.min\_run\_observations} & 2 & Debounce for a new run\\
\texttt{compaction.hot\_days} / \texttt{important\_hot\_days} & 3 / 7 & Detail retention\\
\texttt{compaction.weekly\_rollup\_after\_days} / \texttt{monthly\_\ldots} & 30 / 365 & Roll-up\\
\texttt{evidence.max\_side}, \texttt{quality} & 320, 85 & Crop size and JPEG quality\\
\bottomrule\end{tabularx}\end{table}

Set-up and the main commands (from \texttt{README.md}):
\begin{verbatim}
uv sync                                  # install
uv run pytest                            # 451 tests, coverage floor 75%
uv run aegis-ingest --url http://<phone>:8080/video
uv run python run_pipeline.py --live-session <id>
modal deploy src/aegis/remote/serve.py   # aegis-memory-api
uv run aegis-ask "where are my keys?"
\end{verbatim}

\section{Tools, Frameworks and Technologies}
In addition to the Iteration~1 stack: SAM~3 (default detector), DINOv3
(re-identification), SQLite, the Flutter \texttt{speech\_to\_text} plugin
(on-device speech), and \texttt{matplotlib} and PlantUML for the report figures.
Tools used for writing and coding are disclosed in Section~\ref{sec:tools}.

\section{Demonstration Evidence}
\tbd{Screenshots to be added by the team: Patient App home and answer sheet,
Caregiver App Interior and Track tabs, Rerun view of walk01-2, and a terminal
run of \texttt{run\_pipeline.py}.}

\section{Iteration Outcome}
\label{sec:iter2-outcome}
Objectives (1)--(3) were met as implemented slices. A live capture now reaches
memory on Modal (end-to-end run of 20~November, Section~\ref{sec:e2e}); memory is
bounded (5.7--7.2$\times$ smaller, 100\% retention of last-known states,
Section~\ref{sec:compaction-audit}) and persists across sessions (12/12 instances
still found 1--4 days later on EgoLife, Section~\ref{sec:egolife}); and both apps
are built and tested. Not met: temporal answers (FR5.2, deviation D2),
calibrated abstention, authentication (FR8.1), encryption, the one-hour soak,
latency p95 and any user study. Room linking proved unstable across scale
(TC12). These items form Iteration~3 and~4 (Chapter~\ref{chap:fyp2}).
